\section{Evidence-Allocation Framework and Legal Scope}
\label{app:framework_submission}

The model used in the main text is intentionally spare. Its job
is to justify an ordering rule for scarce forensic attention,
not to model the full internal organization of a cartel. Issues
such as side payments, discipline, and leniency incentives are
treated in the cartel literature
(Asker, 2010; Marshall and Marx, 2012). The narrower
object here is the pre-proof enforcement problem: why a firm
that repeatedly loses, while continuing to appear in procurement
records, can be a useful target for bid-level follow-up.

\subsection{Primitives}
\label{app:framework_primitives}

Consider a procurement environment with repeated tender-items.
In a cartel-targeted tender, one firm is assigned the winning
role and submits the designated bid $b^*$. A losing-role
participant submits a bid $b_\ell>b^*$ and is not meant to win.
The losing role has a cost $c_1>0$. Winning by mistake or by
deviation creates an internal cartel sanction $\kappa>0$,
interpreted broadly as the loss of cartel discipline, future
rents, or capacity-allocation rights. Let $r$ denote the private
gain from deviating into the winner role in a tender where the
firm was assigned to lose.

Firms outside the winning role are of two reduced-form types.
Type $C$ firms are cartel-deployed losing-role participants. Type
$G$ firms are ordinary losing participants: they may be
inefficient, capacity constrained, inexperienced, or simply
unlucky. The type notation is a device for stating the maintained
ranking condition; it is not an empirical classification rule.
The empirical screen conditions on the observed event
$W_i=0$, where $W_i$ is the firm's number of wins, and ranks
firms by participation count $T_i$.

\begin{lemma}[Zero-win losing role]
\label{lem:submission_zero_win}
If the internal sanction from winning in a losing role exceeds
the private deviation gain, $\kappa>r$, a losing-role firm's
equilibrium bid lies above the designated winning bid:
$b_\ell>b^*$. Hence $\Pr(W_i>0\mid C)=0$ over tenders in which
the firm is deployed in the losing role.
\end{lemma}

\begin{proof}
If the losing-role participant submits $b_\ell>b^*$, it loses
and receives the deployment payoff net of the role cost. If it
deviates to a winning bid, it gains $r$ but pays the sanction
$\kappa$. The deviation is unprofitable when $\kappa>r$. The
result fixes the win probability of the losing role, not the
exact numerical bid. Any bid above the designated winner's bid is
consistent with the reduced-form role. 
\end{proof}

Lemma~\ref{lem:submission_zero_win} explains why the zero-win
conditioning set is a coherent place to begin. The conditioning
set is intentionally broad: many genuine firms also never win.
The evidentiary content comes from the second observable,
participation intensity.

\subsection{Participation as a Ranking Statistic}
\label{app:framework_participation}

Within the zero-win set, let participation counts follow a
Poisson approximation:
\[
  T_i \mid \tau_i, q \sim \text{Poisson}(\lambda_q \tau_i),
  \qquad q\in\{C,G\},
\]
where $\tau_i$ is the length of time firm $i$ is observed and
$\lambda_q$ is the type-specific participation rate. The
screening assumption is
\[
  \lambda_C>\lambda_G
  \quad\text{on the conditioning set } W_i=0.
\]
This is the substantive assumption behind the frequent-loser
ranking. Cartel-deployed losing-role participants, if present,
are repeatedly deployed to provide losing participation; ordinary
zero-win firms participate less often.
The condition $\lambda_C>\lambda_G$ is \emph{maintained}, not
derived: it is the behavioral assumption under which repeated
zero-win participation becomes a monotone ranking statistic for
loser-side exposure. The empirical sections then ask whether
the ranking implied by this assumption has validation content
against adjudication-adjacent labels.
The type notation is a device for stating the maintained ranking
condition; it is not an empirical classification rule.
% CO-AUTHOR EDIT [v21 Major 4]: disclose the observational equivalence.
Two distinct behavioral processes generate $\lambda_C>\lambda_G$
on the conditioning set: deployment of cover bidders during a
cartel's operation, and sincere persistent bidding by low-cost
firms that rarely win. These are observationally close, and
conditioning on $W_i=0$ favors the latter through survivorship.
The deployment-timing audit in
\ref{app:deployment_timing_submission} finds a partial
separation---cobidder participation tilts toward the years of
cartel contact---but does not exclude sincere persistence. The
ranking is therefore used as a predictive triage statistic, not
as a causal test of cartel role.

\begin{proposition}[Monotone loser-side suspicion]
\label{prop:submission_monotone}
Under the Poisson approximation and $\lambda_C>\lambda_G$, the
posterior probability $\Pr(C\mid T_i=t,W_i=0,\tau_i)$ is
non-decreasing in $t$. Consequently, $T_i$ and
$\log(1+T_i)$ are rank-equivalent screening statistics for
loser-side cartel-adjacency exposure within the zero-win
stratum.
\end{proposition}

\begin{proof}
By Bayes' rule, the posterior odds of type $C$ against type $G$
are proportional to the likelihood ratio
\[
  \frac{\Pr(T_i=t\mid C,\tau_i)}
       {\Pr(T_i=t\mid G,\tau_i)}
  =
  \exp[-(\lambda_C-\lambda_G)\tau_i]
  \left(\frac{\lambda_C}{\lambda_G}\right)^t .
\]
Because $\lambda_C>\lambda_G$, the likelihood ratio is increasing
in $t$. This is the monotone-likelihood-ratio property of the
Poisson family (Karlin and Rubin, 1956).%
\footnote{Karlin, S. and H.~Rubin (1956). The theory of decision
procedures for distributions with monotone likelihood ratio.
\textit{Annals of Mathematical Statistics} 27(2), 272--299.}
The posterior increases
with participation count, and the log transformation
$\log(1+T_i)$ preserves the ranking.
\end{proof}

The empirical binary rule used in the main text is an operational
coarsening of Proposition~\ref{prop:submission_monotone}. It flags firms above the
median plus $1.5$ times the interquartile range of zero-win
participation. The threshold is not a structural parameter. It
is an auditable way to turn a monotone ranking into an
operational list. Agencies could instead use the continuous
score, a capacity-constrained top-$k$ list, or a threshold
calibrated to their investigative budget.
The posterior language belongs to the reduced-form screening
model. The empirical implementation uses the statistic as a
ranking device for forensic priority, not as an estimated
probability of legal membership.

\subsection{Evidentiary Scope}
\label{app:framework_scope}

The proposition delivers a ranking result. Even if the model is
correct, a high value of
$\log(1+T_i)$ raises the posterior probability of loser-side
cartel-adjacency exposure only relative to other zero-win firms.
It does not
identify an agreement, assign cartel membership, prove that any
specific bid was cover, or establish a damages base. Those
claims require evidence from the richer record: bid values,
communications, leniency material, repeated winner rotation, or
other case-specific facts.

This boundary is the legal reason for the staged architecture in
the main text. The award layer can allocate attention because it
ranks firms on a statistic that is cheap, observable, and
monotone under a transparent behavioral assumption. The bid layer
and case record then do the evidentiary work that liability
requires, because liability turns on conduct and agreement, not
on suspicious participation alone.

The framework does not imply that all zero-win firms occupy a
cartel losing role, that frequent-loser status is legal
membership, or that any specific bid is a cover bid. It supplies
only the ranking logic for triage.

% CO-AUTHOR EDIT [v21 minor]: the former "Proposition 2" was a definitional
% restatement, not a derived result; demote to a scope remark.
\paragraph{Scope of the ranking} The monotonicity result is a
\emph{ranking} rationale, not a structural model. It orders
expected investigative value within the zero-win stratum; it says
nothing about liability or damages, which require different
predicates---proof of agreement, conduct in particular tenders,
and an economic measure of price injury. The statistic carries
participation and win information only and omits the facts those
predicates need, so the same object can be useful for allocating
forensic attention and insufficient for legal adjudication. This
is a definitional boundary on what the ranking claims, not a
separate theorem.

% CO-AUTHOR EDIT [v21 conceptual framing]: proof sketches for the body's
% decision-theoretic frame (Section ref{sec:triage-model}). Reuses the
% maintained MLRP above; does not re-derive it.
\subsection{Triage Model: Proof Sketches}
\label{app:triage-proofs}

These sketches support the decision-theoretic frame of
\S\ref{sec:triage-model}. They use the maintained within-stratum
MLRP of \S\ref{app:framework_participation}
($\lambda(s\mid e)=f(s\mid\theta=1,e)/f(s\mid\theta=0,e)$
nondecreasing in $s$); no claim about how $\theta=1$ firms bid is
invoked.

\paragraph{Proposition~\ref{prop:index} (index rule).}
By Bayes' rule with within-stratum prior $\pi(e)$,
$\rho(s,e)=[1+\tfrac{1-\pi(e)}{\pi(e)}\lambda(s\mid e)^{-1}]^{-1}$,
a strictly increasing transform of $\lambda(s\mid e)$. The yield
\eqref{eq:agency-yield} is additively separable across firms under
the cardinality constraint $|A|\le K$ (a linear objective over a
uniform matroid), so the greedy rule selecting the $K$ largest
$\rho(s_i,e_i)$ is optimal; equivalently, any set omitting a
higher-$\rho$ firm for a lower-$\rho$ one is weakly improved by the
swap. Under MLRP, $\rho(s,e)$ is nondecreasing in $s$ within each
stratum, so ordering by $\rho$ coincides with ordering by $s$ and
the optimum is a per-stratum cutoff on the award rank.
\hfill$\square$

\paragraph{Proposition~\ref{prop:incremental} (incremental value over exposure).}
\emph{Only if.} If $\theta\perp s\mid e$ then
$\rho(s,e)=\rho_e(e)$ pointwise, the two scores agree as
functions, and both policies escalate the same $K$ firms (up to
ties): the yields coincide. \emph{If.} Suppose
$\theta\not\perp s\mid e$. Conditional dependence \emph{together
with} MLRP yields a stratum $e^\star$ in which
$\lambda(s\mid e^\star)$ is strictly increasing over a
positive-measure set, forcing first-order stochastic dominance of
$s\mid(\theta=1,e^\star)$ over $s\mid(\theta=0,e^\star)$, hence
$\mathrm{AUC}(s,\theta\mid e^\star)>\tfrac12$ strictly. (This MLRP
step closes the equivalence: absent MLRP, distinct conditional
laws can give $\mathrm{AUC}=\tfrac12$---e.g.\ a symmetric
spread-only difference---so dependence alone does not imply rank
discrimination above one half.) The exposure-only score is
constant on the stratum and breaks ties
$\theta$-uninformatively, so its per-slot expected target content
equals the stratum mean $\pi(e^\star)$; with the budget escalating
some but not all of the stratum, the $s$-policy escalates the
highest-$s$ firms, whose per-slot content strictly exceeds
$\pi(e^\star)$ by the rank inequality, so $V$ is strictly larger.
Equivalently, the $\rho$-optimal set of
Proposition~\ref{prop:index} weakly dominates the $\rho_e$-set,
strictly whenever the latter is not $\rho$-optimal---which the
binding-budget-within-an-informative-stratum condition guarantees.
\hfill$\square$

\noindent\emph{Caveat (stated, not buried).} The strict inequality
needs the budget to bind \emph{within} an informative stratum and
the exposure-only tie-break to be $\theta$-uninformative. Under a
slack budget, or a tie-break that uses $\theta$-correlated
information, the selections can coincide even under conditional
dependence, leaving only a weak ($\ge$) relation; the ``iff''
holds under the stated binding-budget restriction, not globally.

\paragraph{Remark~\ref{rem:oos} (out-of-sample crossing): toy Gaussian.}
We make the bias--variance reversal precise in a stylized model and
fence its scope; it is not distribution-free. Each firm carries a
true score $\mu_i$ (its policy posterior); the agency acts on an
estimated score $\widehat m_i=\mu_i+\varepsilon_i$,
$\varepsilon_i\sim\mathcal N(0,\sigma^2)$, and escalates the
top-$K$ by $\widehat m$. The sequential gate on $s$ has
estimation-noise variance $\sigma_s^2$ (one low-variance count);
the joint score on $(s,b)$ has $\sigma_J^2>\sigma_s^2$ (many
estimated moments). In sample the joint score is weakly more
informative, with informational gain $G_b\ge0$ (data processing).
Expected top-$K$ yield is governed by the standardized separation
$d(\sigma)=\Delta\mu/\sqrt{\nu+\sigma^2}$ (with $\Delta\mu$ the
true mean-score gap and $\nu$ the noise-free dispersion): it is
strictly increasing in $d$ and $d$ is strictly decreasing in
$\sigma^2$, so expected yield $Y(\sigma^2)$ is nonincreasing in
scoring noise, from $Y(0)$ (in-sample optimum) to $Y(\infty)=K\pi$
(base-rate floor). Writing $Y_J=Y_J(0)-L_J(\sigma_J^2)$ and
$Y_s=Y_s(0)-L_s(\sigma_s^2)$ with increasing noise losses
$L_\cdot$, the gate wins out of sample,
$Y_s(\sigma_s^2)>Y_J(\sigma_J^2)$, exactly when
$L_J(\sigma_J^2)-L_s(\sigma_s^2)>G_b$. \emph{Existence is not
automatic}: since $L_J\le Y_J(0)-K\pi$, a finite crossing requires
$G_b<\sup_\sigma L_J-L_s(\sigma_s^2)$---if $b$'s in-sample gain is
large, no finite noise reverses the order. The most fragile
idealization is independence of $\varepsilon_i$ across firms:
scores fit on a shared training sample have correlated coefficient
error, which is precisely the out-of-sample degradation a temporal
holdout proxies. Hence we keep this a Remark; the temporal-holdout pattern
reported in \S\ref{sec:triage-model} (Remark~\ref{rem:oos})---the
sequential gate matching or edging the joint score within the
gating budget range---is consistent with the high-relative-noise
regime, not a verification of the toy model, and the advantage is
confined to that range.
\hfill$\square$

\paragraph{Corollary~\ref{cor:separation} (attention vs.\ liability).}
The triage index is a transform of $s$. If
$s\perp\mathbf{1}[i\in D]$, the conditional law of $s$ is identical
for principals and non-principals, so the likelihood ratio of $s$
for membership in $D$ is identically one and
$\mathrm{AUC}(s,\mathbf{1}[i\in D])=\tfrac12$ (independence
delivers this directly; no MLRP needed for the liability label).
The same $s$ satisfies $\mathrm{AUC}(s,\theta\mid e)>\tfrac12$ by
Proposition~\ref{prop:incremental}, so the index can be
informative for the escalation target while remaining near-random
for principal status. The empirical near-random value is evidence
consistent with the role-allocation premise, not an estimation
failure. \hfill$\square$
